\section{Introduction}
\label{sec:intro}

Large language models are increasingly deployed in production systems that process continuously evolving input streams---from customer support queries to news summarization pipelines. When the input distribution shifts, model predictions can silently degrade without any explicit error signal~\cite{khaki2023uncovering}. Detecting such distribution drift early is critical for maintaining system reliability.

Standard drift detectors monitor low-order statistics of embedding distributions: centroid shift tracks the mean, covariance shift tracks the second moment, and Maximum Mean Discrepancy (MMD) measures distributional distance in a reproducing kernel Hilbert space~\cite{khaki2023uncovering, sodar2023detecting}. These methods are effective when drift manifests as shifts in location or spread, but they may miss subtler forms of drift that reorganize the \emph{geometric structure} of the embedding manifold without changing its center of mass. For example, if user queries shift from a diverse cluster to a ring-like distribution around a fixed topic---avoiding certain sub-topics while converging on others---the centroid remains stable while the topology changes fundamentally.

Topological Data Analysis (TDA), and persistent homology in particular, captures multi-scale geometric features of point clouds that are invariant to coordinate transformations~\cite{tralie2019ripser, tauzin2020giotto}. Recent work has applied TDA to characterize LLM representations~\cite{gardinazzi2024persistent}, detect LLM-generated text~\cite{wei2025shortphd}, and monitor concept drift in image streams~\cite{basterrech2024unsupervised}. However, no prior work systematically compares TDA-based drift detectors against classical baselines on real LLM embedding streams across multiple drift types.

We address this gap with the first controlled comparative study of TDA versus classical drift detection on sentence embedding streams. Using \agnews embeddings from \minilm (\texttt{all-MiniLM-L6-v2}, 384 dimensions), we construct five drift scenarios---abrupt topic drift, gradual topic drift, geometric drift, style shift, and a no-drift control---and evaluate nine drift detectors spanning both statistical and TDA families. Our main findings are:

\begin{itemize}[leftmargin=*,itemsep=0pt,topsep=0pt]
    \item We conduct the first systematic comparison of TDA-based drift detectors (persistent entropy, Wasserstein distance, PHD, Betti curves) against classical baselines (centroid shift, covariance shift, MMD, kNN) on real LLM embedding streams.
    \item We show that classical statistical methods achieve perfect detection (AUC\,=\,1.0) on natural topic drift in \agnews embeddings, while the best TDA method (Wasserstein H0) reaches AUC\,=\,0.93 on abrupt drift.
    \item We design a synthetic topology-only drift experiment that isolates TDA's genuine advantage: H1 persistent entropy achieves $12.3\times$ separability on centroid-preserving loop drift, where centroid shift fails entirely ($0.8\times$).
    \item We provide practical guidelines for combining TDA and statistical detectors in a complementary monitoring pipeline, with runtime analysis showing TDA adds only $\sim$37ms overhead per window on CPU.
\end{itemize}

The rest of this paper is organized as follows. \Secref{sec:related} reviews related work on drift detection and TDA for NLP. \Secref{sec:method} describes our methodology, including drift scenarios, detectors, and evaluation protocol. \Secref{sec:results} presents experimental results on both real and synthetic data. \Secref{sec:discussion} discusses implications and limitations. \Secref{sec:conclusion} concludes with practical recommendations and future directions.
